\chapter{Програмын үндсэн кодын блокууд}
\label{AppendixC}
\pagecolor{white}

Энэ хавсралтад 2-р бүлэгт тайлбарласан программ хангамжийн загварчлал ба хөгжүүлэлтэд ашигласан үндсэн кодын блокуудыг бүрэн эхээр нь нэгтгэн оруулав. Бүлэгт зөвхөн ишлэлээр хийсвэр зам нэмж тайлбарлан, бодит код энэ хавсралтад байрлана.

%-------------------------------------------------------------------------------

\section*{В.1 Flutter --- AppTheme класс (Material Design 3)}

\begin{lstlisting}[language={}, caption={lib/theme/app\_theme.dart --- Дулаан говийн өнгөн схем}, label={lst:apptheme}]
class AppTheme {
  static const Color primary       = Color(0xFFEDE7DA);
  static const Color secondary     = Color(0xFFD4700F);
  static const Color accent        = Color(0xFFB85C08);
  static const Color textPrimary   = Color(0xFF1A1208);

  static ThemeData get lightTheme => ThemeData(
    useMaterial3: true,
    fontFamily: 'Montserrat',
    brightness: Brightness.light,
    scaffoldBackgroundColor: primary,
    colorScheme: const ColorScheme.light(
      primary:   secondary,
      secondary: accent,
    ),
  );
}
\end{lstlisting}

%-------------------------------------------------------------------------------

\section*{В.2 Flutter --- BottomNavigationBar (4 tab навигаци)}

\begin{lstlisting}[language={}, caption={lib/screens/main\_screen.dart --- үндсэн навигаци}, label={lst:bottomnav}]
class _MainScreenState extends State<MainScreen> {
  int _selectedIndex = 0;
  final List<Widget> _screens = const [
    HomeScreen(), BookingsScreen(),
    ChatScreen(), ProfileScreen(),
  ];

  @override
  Widget build(BuildContext context) {
    return Scaffold(
      body: _screens[_selectedIndex],
      bottomNavigationBar: BottomNavigationBar(
        currentIndex: _selectedIndex,
        onTap: (index) => setState(() => _selectedIndex = index),
        items: const [
          BottomNavigationBarItem(icon: Icon(Icons.home), label: 'Home'),
          BottomNavigationBarItem(icon: Icon(Icons.calendar_today), label: 'Bookings'),
          BottomNavigationBarItem(icon: Icon(Icons.chat), label: 'Chat'),
          BottomNavigationBarItem(icon: Icon(Icons.person), label: 'Profile'),
        ],
      ),
    );
  }
}
\end{lstlisting}

%-------------------------------------------------------------------------------

\section*{В.3 Flutter --- Firebase Email Authentication}

\begin{lstlisting}[language={}, caption={lib/services/auth\_service.dart --- бүртгэлийн логик}, label={lst:auth}]
Future<void> registerWithEmail(
    String email, String password, String name) async {
  final credential = await FirebaseAuth.instance
      .createUserWithEmailAndPassword(
        email: email, password: password);
  final uid = credential.user!.uid;

  await http.post(
    Uri.parse(AppConfig.authEndpoint + '/send-code'),
    headers: {'Content-Type': 'application/json'},
    body: jsonEncode({'email': email, 'uid': uid}),
  );
}
\end{lstlisting}

%-------------------------------------------------------------------------------

\section*{В.4 Flutter --- Цагийн хуваарь ачаалах логик}

\begin{lstlisting}[language={}, caption={lib/screens/booking\_screen.dart --- захиалагдсан цагийн үеийг тооцоолох}, label={lst:slots}]
Future<void> _loadSlots(String venueId, String date) async {
  final bookSnap = await FirebaseFirestore.instance
      .collection('bookings')
      .where('venueId', isEqualTo: venueId)
      .where('date', isEqualTo: date)
      .where('status', whereIn: ['upcoming', 'active'])
      .get();

  final weekday = DateTime.parse(date).weekday % 7;
  final fixedSnap = await FirebaseFirestore.instance
      .collection('fixed_bookings')
      .where('venueId', isEqualTo: venueId)
      .where('dayOfWeek', isEqualTo: weekday)
      .where('isActive', isEqualTo: true)
      .get();

  final bookedSlots = {
    ...bookSnap.docs.map((d) => d['timeSlot'] as String),
    ...fixedSnap.docs.map((d) => d['timeSlot'] as String),
  };
  setState(() => _bookedSlots = bookedSlots);
}
\end{lstlisting}

%-------------------------------------------------------------------------------

\section*{В.5 Firebase Cloud Functions --- Claude Haiku AI chat endpoint}

\begin{lstlisting}[language={}, caption={functions/index.js --- Anthropic Claude Haiku-тай холбогдох onCall функц}, label={lst:claude}]
exports.chat = functions.https.onCall(async (data, context) => {
  if (!context.auth) throw new functions.https.HttpsError(
      'unauthenticated', 'Authentication required');
  const client = new Anthropic({ apiKey: functions.config().anthropic.key });
  const message = await client.messages.create({
    model: 'claude-haiku-4-5', max_tokens: 512,
    messages: [{ role: 'user', content: data.message }],
    system: 'You are a sports hall booking assistant for Dalanzadgad city.',
  });
  return { reply: message.content[0].text };
});
\end{lstlisting}

%-------------------------------------------------------------------------------

\section*{В.6 Express.js --- үндсэн route тохиргоо (server.js)}

\begin{lstlisting}[language={}, caption={backend/server.js --- Express middleware ба route модулиуд}, label={lst:server}]
const express = require('express');
const app     = express();
app.use(express.json());
app.use('/api/auth',     require('./routes/auth'));
app.use('/api/chat',     require('./routes/chat'));
app.use('/api/users',    require('./routes/users'));
app.use('/api/bookings', require('./routes/bookings'));
app.use('/api/schedule', require('./routes/schedule'));
app.listen(3000, () => console.log('Server running on port 3000'));
\end{lstlisting}

%-------------------------------------------------------------------------------

\section*{В.7 Backend --- Cron ажил (15 минут тутамд захиалга шинэчлэх)}

\begin{lstlisting}[language={}, caption={backend/cron/update\_bookings.js --- node-cron schedule}, label={lst:cron}]
cron.schedule('*/15 * * * *', async () => {
  const now   = new Date();
  const today = now.toISOString().split('T')[0];
  const nowTime = now.getHours().toString().padStart(2,'0') + ':'
                + now.getMinutes().toString().padStart(2,'0');
  const snapshot = await admin.firestore()
    .collection('bookings')
    .where('status', '==', 'upcoming')
    .where('dateKey', '<=', today).get();
  const batch = admin.firestore().batch();
  snapshot.docs.forEach(doc => {
    const endTime = doc.data().timeSlotEnd;
    if (doc.data().dateKey < today ||
        (doc.data().dateKey === today && endTime <= nowTime)) {
      batch.update(doc.ref, { status: 'completed' });
    }
  });
  await batch.commit();
});
\end{lstlisting}
